\documentclass[10pt,a4paper]{amsart}
\usepackage{amsmath,amssymb,amsthm,graphicx}
\usepackage{verbatim}
\newenvironment{reference}{\comment}{\endcomment}
%\newenvironment{reference}{}{}
\newenvironment{slogan}{\comment}{\endcomment}
\newenvironment{history}{\comment}{\endcomment}

% For commutative diagrams we use Xy-pic
\usepackage[all]{xy}

% We use 2cell for 2-commutative diagrams.
\xyoption{2cell}
% Theorem environments.
%
\theoremstyle{plain}
\newtheorem{theorem}[subsection]{Theorem}
\newtheorem{proposition}[subsection]{Proposition}
\newtheorem{lemma}[subsection]{Lemma}

\theoremstyle{definition}
\newtheorem{definition}[subsection]{Definition}
\newtheorem{example}[subsection]{Example}
\newtheorem{exercise}[subsection]{Exercise}
\newtheorem{situation}[subsection]{Situation}

\theoremstyle{remark}
\newtheorem{remark}[subsection]{Remark}
\newtheorem{remarks}[subsection]{Remarks}

\numberwithin{equation}{subsection}

% Macros
%
\def\lim{\mathop{\mathrm{lim}}\nolimits}
\def\colim{\mathop{\mathrm{colim}}\nolimits}
\def\Spec{\mathop{\mathrm{Spec}}}
\def\Hom{\mathop{\mathrm{Hom}}\nolimits}
\def\Ext{\mathop{\mathrm{Ext}}\nolimits}
\def\SheafHom{\mathop{\mathcal{H}\!\mathit{om}}\nolimits}
\def\SheafExt{\mathop{\mathcal{E}\!\mathit{xt}}\nolimits}
\def\Sch{\mathit{Sch}}
\def\Mor{\mathop{\mathrm{Mor}}\nolimits}
\def\Ob{\mathop{\mathrm{Ob}}\nolimits}
\def\Sh{\mathop{\mathit{Sh}}\nolimits}
\def\NL{\mathop{N\!L}\nolimits}
\def\CH{\mathop{\mathrm{CH}}\nolimits}
\def\proetale{{pro\text{-}\acute{e}tale}}
\def\etale{{\acute{e}tale}}
\def\QCoh{\mathit{QCoh}}
\def\Ker{\mathop{\mathrm{Ker}}}
\def\Im{\mathop{\mathrm{Im}}}
\def\Coker{\mathop{\mathrm{Coker}}}
\def\Coim{\mathop{\mathrm{Coim}}}

% Boxtimes
%
\DeclareMathSymbol{\boxtimes}{\mathbin}{AMSa}{"02}

%
% Macros for moduli stacks/spaces
%
\def\QCohstack{\mathcal{QC}\!\mathit{oh}}
\def\Cohstack{\mathcal{C}\!\mathit{oh}}
\def\Spacesstack{\mathcal{S}\!\mathit{paces}}
\def\Quotfunctor{\mathrm{Quot}}
\def\Hilbfunctor{\mathrm{Hilb}}
\def\Curvesstack{\mathcal{C}\!\mathit{urves}}
\def\Polarizedstack{\mathcal{P}\!\mathit{olarized}}
\def\Complexesstack{\mathcal{C}\!\mathit{omplexes}}
% \Pic is the operator that assigns to X its picard group, usage \Pic(X)
% \Picardstack_{X/B} denotes the Picard stack of X over B
% \Picardfunctor_{X/B} denotes the Picard functor of X over B
\def\Pic{\mathop{\mathrm{Pic}}\nolimits}
\def\Picardstack{\mathcal{P}\!\mathit{ic}}
\def\Picardfunctor{\mathrm{Pic}}
\def\Deformationcategory{\mathcal{D}\!\mathit{ef}}
\usepackage{cleveref}
\crefname{originallemma}{original lemma}{original lemmas}
\Crefname{originallemma}{Original lemma}{Original lemmas}
\makeatletter
\newlabel{lemma-kernel-finite-locally-constant}{{03RX}{source}{}{sourcepointer-1}{}}
\newlabel{lemma-kernel-finite-locally-constant@cref}{{[originallemma][0][]03RX}{[1][0][]source}{}{}{}}
\makeatother
\crefname{originallemma}{original lemma}{original lemmas}
\Crefname{originallemma}{Original lemma}{Original lemmas}
\makeatletter
\newlabel{lemma-support-constructible}{{09YS}{source}{}{sourcepointer-2}{}}
\newlabel{lemma-support-constructible@cref}{{[originallemma][0][]09YS}{[1][0][]source}{}{}{}}
\makeatother
\crefname{originallemma}{original lemma}{original lemmas}
\Crefname{originallemma}{Original lemma}{Original lemmas}
\makeatletter
\newlabel{properties-lemma-quasi-compact-quasi-separated-spectral}{{094L}{source}{}{sourcepointer-3}{}}
\newlabel{properties-lemma-quasi-compact-quasi-separated-spectral@cref}{{[originallemma][0][]094L}{[1][0][]source}{}{}{}}
\makeatother
\crefname{originallemma}{original lemma}{original lemmas}
\Crefname{originallemma}{Original lemma}{Original lemmas}
\makeatletter
\newlabel{topology-lemma-constructible-hausdorff-quasi-compact}{{0901}{source}{}{sourcepointer-4}{}}
\newlabel{topology-lemma-constructible-hausdorff-quasi-compact@cref}{{[originallemma][0][]0901}{[1][0][]source}{}{}{}}
\makeatother
\title{A filtered colimit of constructible sheaves with constructible colimit is essentially constant}
\author{The Stacks Project Authors}
\date{stacks-project:etale-cohomology.tex:local-snapshot-20261011; GNU FDL 1.2 or later}
\begin{document}
\maketitle
\noindent\textit{Editorial note.} All delivered mathematical text is the author's own, copied from the pinned Stacks Project chapter file etale-cohomology.tex; no statement or proof step was rewritten, reordered or removed.\par\medskip
\noindent\textit{Editorial note.} The proof is the author's complete argument for all three parts; the parenthetical that the constructibility of one locus is left to the reader comes with the author's explicit hint naming the result to use, and it is preserved verbatim in the counted span.\par\medskip
\noindent\textit{Editorial note.} Every labelled reference occurring inside the counted spans is carried in assembly.external\_references as a non-load-bearing navigation pointer (same-chapter ones at their definition site, cross-chapter ones at the invocation inside the counted span), so no load-bearing mathematics is replaced by a stub. No bibliography citation occurs inside the counted spans.\par\medskip
\noindent\textit{Editorial note.} The checked chapter list for same-name collisions is the set of stacks chapters staged in this repository: chow, cohomology, divisors, etale-cohomology, flat, more-algebra, more-morphisms, morphisms, spaces-more-morphisms, varieties together with this card's chapter etale-cohomology.\par\medskip
\noindent\textit{Editorial note.} The preamble is the chapter's own native preamble excerpt already used for this source: the Xy-pic setup (native preamble lines 14-18) and the native theorem declarations with the mathematical macros (native preamble lines 154-217).\par\medskip
\noindent\textit{Allowed standard prerequisites (editorial).}\begin{itemize}
\item Constructible sheaves are closed under finite limits and colimits. Same chapter, Lemma lemma-kernel-finite-locally-constant. Supplies that the sheaves appearing in the constructions are constructible and that constructible sheaves of modules form an abelian category.
\item Constructibility of the support of a map of constructible sheaves. Same chapter, Lemma lemma-support-constructible. Makes the loci where the maps on stalks fail to be isomorphisms constructible.
\item The constructible topology on a quasi-compact quasi-separated scheme is quasi-compact. Properties, Lemma properties-lemma-quasi-compact-quasi-separated-spectral and Topology, Lemma topology-lemma-constructible-hausdorff-quasi-compact. Converts pointwise eventual emptiness of a constructible subset into emptiness.
\end{itemize}\medskip
\ifdefined\hypertarget\hypertarget{sourcepointer-1}{}\fi
\noindent\textit{Source reference.} Original lemma 03RX: sources/stacks-etale-cohomology/pinned-original-source.tex, source lines 10050--10050. Navigation pointer invoked at proof line 11860; the counted span contains the pointer only and the cited result is not reproduced inside it, so the pointer is registered as a non-load-bearing navigation stub in the established form for Stacks material. No bibliography citation is involved.\par\medskip
\ifdefined\hypertarget\hypertarget{sourcepointer-2}{}\fi
\noindent\textit{Source reference.} Original lemma 09YS: sources/stacks-etale-cohomology/pinned-original-source.tex, source lines 11867--11867. Navigation pointer invoked at proof line 11929; the counted span contains the pointer only and the cited result is not reproduced inside it, so the pointer is registered as a non-load-bearing navigation stub in the established form for Stacks material. No bibliography citation is involved.\par\medskip
\ifdefined\hypertarget\hypertarget{sourcepointer-3}{}\fi
\noindent\textit{Source reference.} Original lemma 094L: sources/stacks-etale-cohomology/pinned-original-source.tex, source lines 11934--11934. Navigation pointer invoked at proof line 11934; the counted span contains the pointer only and the cited result is not reproduced inside it, so the pointer is registered as a non-load-bearing navigation stub in the established form for Stacks material. No bibliography citation is involved.\par\medskip
\ifdefined\hypertarget\hypertarget{sourcepointer-4}{}\fi
\noindent\textit{Source reference.} Original lemma 0901: sources/stacks-etale-cohomology/pinned-original-source.tex, source lines 11936--11936. Navigation pointer invoked at proof line 11936; the counted span contains the pointer only and the cited result is not reproduced inside it, so the pointer is registered as a non-load-bearing navigation stub in the established form for Stacks material. No bibliography citation is involved.\par\medskip
\begin{lemma}
\label{lemma-colimit-constructible}
Let $X$ be a quasi-compact and quasi-separated scheme. Let
$\mathcal{F} = \colim_{i \in I} \mathcal{F}_i$ be a filtered
colimit of sheaves of sets, abelian sheaves, or sheaves of modules.
\begin{enumerate}
\item If $\mathcal{F}$ and $\mathcal{F}_i$ are constructible sheaves of
sets, then the ind-object $\mathcal{F}_i$ is essentially constant with
value $\mathcal{F}$.
\item If $\mathcal{F}$ and $\mathcal{F}_i$ are constructible sheaves of
abelian groups, then the ind-object $\mathcal{F}_i$ is essentially constant
with value $\mathcal{F}$.
\item Let $\Lambda$ be a Noetherian ring.
If $\mathcal{F}$ and $\mathcal{F}_i$ are constructible sheaves of
$\Lambda$-modules, then the ind-object $\mathcal{F}_i$ is essentially constant
with value $\mathcal{F}$.
\end{enumerate}
\end{lemma}

\begin{proof}
Proof of (1). We will use without further mention that finite limits
and colimits of constructible sheaves are constructible
(Lemma \ref{lemma-kernel-finite-locally-constant}).
For each $i$ let $T_i \subset X$ be the set of points $x \in X$
where $\mathcal{F}_{i, \overline{x}} \to \mathcal{F}_{\overline{x}}$
is not surjective. Because $\mathcal{F}_i$ and $\mathcal{F}$ are
constructible $T_i$ is a constructible subset of $X$
(Lemma \ref{lemma-support-constructible}).
Since the stalks of $\mathcal{F}$ are finite
and since $\mathcal{F} = \colim_{i \in I} \mathcal{F}_i$ we see
that for all $x \in X$ we have $x \not \in T_i$ for $i$ large enough.
Since $X$ is a spectral space by Properties, Lemma
\ref{properties-lemma-quasi-compact-quasi-separated-spectral}
the constructible topology on $X$ is quasi-compact by
Topology, Lemma \ref{topology-lemma-constructible-hausdorff-quasi-compact}.
Thus $T_i = \emptyset$ for $i$ large enough. Thus
$\mathcal{F}_i \to \mathcal{F}$ is surjective for $i$ large enough.
Assume now that $\mathcal{F}_i \to \mathcal{F}$ is surjective for all $i$.
Choose $i \in I$. For $i' \geq i$ denote $S_{i'} \subset X$ the set of
points $x$ such that the number of elements in
$\Im(\mathcal{F}_{i, \overline{x}} \to \mathcal{F}_{\overline{x}})$
is not equal to the number of elements in
$\Im(\mathcal{F}_{i, \overline{x}} \to \mathcal{F}_{i', \overline{x}})$.
Because $\mathcal{F}_i$, $\mathcal{F}_{i'}$ and $\mathcal{F}$ are
constructible $S_{i'}$ is a constructible subset of $X$
(details omitted; hint: use Lemma \ref{lemma-support-constructible}).
Since the stalks of $\mathcal{F}_i$ and $\mathcal{F}$
are finite and since $\mathcal{F} = \colim_{i' \geq i} \mathcal{F}_{i'}$
we see that for all $x \in X$ we have $x \not \in S_{i'}$ for $i'$
large enough. By the same argument as above we can find a large $i'$ such
that $S_{i'} = \emptyset$. Thus $\mathcal{F}_i \to \mathcal{F}_{i'}$
factors through $\mathcal{F}$ as desired.

\medskip\noindent
Proof of (2). Observe that a constructible abelian sheaf is a constructible
sheaf of sets. Thus case (2) follows from (1).

\medskip\noindent
Proof of (3). We will use without further mention that the category of
constructible sheaves of $\Lambda$-modules is abelian
(Lemma \ref{lemma-kernel-finite-locally-constant}).
For each $i$ let $\mathcal{Q}_i$ be the cokernel of the map
$\mathcal{F}_i \to \mathcal{F}$. The support $T_i$ of $\mathcal{Q}_i$
is a constructible subset of $X$ as $\mathcal{Q}_i$ is constructible
(Lemma \ref{lemma-support-constructible}).
Since the stalks of $\mathcal{F}$ are finite $\Lambda$-modules
and since $\mathcal{F} = \colim_{i \in I} \mathcal{F}_i$ we see
that for all $x \in X$ we have $x \not \in T_i$ for $i$ large enough.
Since $X$ is a spectral space by Properties, Lemma
\ref{properties-lemma-quasi-compact-quasi-separated-spectral}
the constructible topology on $X$ is quasi-compact by
Topology, Lemma \ref{topology-lemma-constructible-hausdorff-quasi-compact}.
Thus $T_i = \emptyset$ for $i$ large enough. This proves the first
assertion. For the second, assume now that
$\mathcal{F}_i \to \mathcal{F}$ is surjective for all $i$.
Choose $i \in I$. For $i' \geq i$ denote $\mathcal{K}_{i'}$ the
image of $\Ker(\mathcal{F}_i \to \mathcal{F})$ in $\mathcal{F}_{i'}$.
The support $S_{i'}$ of $\mathcal{K}_{i'}$
is a constructible subset of $X$ as $\mathcal{K}_{i'}$ is constructible.
Since the stalks of $\Ker(\mathcal{F}_i \to \mathcal{F})$
are finite $\Lambda$-modules and since
$\mathcal{F} = \colim_{i' \geq i} \mathcal{F}_{i'}$ we see
that for all $x \in X$ we have $x \not \in S_{i'}$ for $i'$ large enough.
By the same argument as above we can find a large $i'$ such
that $S_{i'} = \emptyset$. Thus $\mathcal{F}_i \to \mathcal{F}_{i'}$
factors through $\mathcal{F}$ as desired.
\end{proof}

\end{document}
